% Stage 8 restored fragment — 2.0.1 Radiation as a source-to-field operator (archive lines 111-549)
\subsection{Radiation as a source-to-field operator}
\label{sec:ch201}

A network analyzer $|S_{11}|$, a spherical near-field scan on $\Gamma_R$, and a far-zone pattern cut are three readings on one ordered map: impressed data, Maxwell closure, outgoing export, finite observation \cite{maxwell1873,stratton1941}. Misplacing a measurement on that chain---reporting reactive storage as radiated power---remains the common compact-range failure mode \cite{poynting1884,hansen1988}.

Section~\ref{sec:ch20} posed radiation as a question of accessibility rather than of
field existence alone. The laboratory scene below makes that distinction operational by naming
the maps that connect impressed currents to Maxwell field pairs and, through outgoing
selection, to far-zone observables \cite{stratton1941}. The distinction between field
existence and field $\mathcal{E}_{\mathrm{acc}}^{(N)}$ is central to radiation theory \cite{stratton1941}.
Maxwell's equations determine which fields may arise from a specified source distribution,
whereas the radiation map $\mathcal{R}_\omega$ determines which portion of those fields
survives as observable outgoing content \cite{maxwell1873}. From an antenna
engineer's perspective, $\mathcal{S}_\omega$ answers what the structure \emph{stores and
circulates} locally, while $\mathcal{R}_\omega$ answers what the range may \emph{meter and
pattern} after the reactive shell is left behind \cite{harrington2001}.

A network analyzer $|S_{11}|$, a spherical near-field scan on $\Gamma_R$, and a far-zone pattern cut are three readings on one ordered map: impressed data, Maxwell field specification, outgoing export, finite observation \cite{maxwell1873}.
Each instrument samples a different stage of that map, and confusion between stages---for
example, reporting near-field reactive content as radiated power---remains the common
compact-range laboratory error \cite{poynting1884}.

Laboratory measurements presuppose an excitation model even when it is left implicit: a port
drive, a probe sample, or a termination fixes how energy is partitioned between storage and
escape \cite{stratton1941}. Impressed data from Section~\ref{sec:ch01.1.2} map through Maxwell
field specification to \((\E,\H)\), and only then is the outward-propagating fraction
identified as the carrier of far-zone power and angular flux
(Eq.~\eqref{eq:ch1.maxwell-local}; Sections~\ref{sec:ch01.1.4}--\ref{sec:ch01.1.5}). At fixed
$\omega$ the maps are the harmonic solution operator $\mathcal{S}_\omega$ and radiation
operator $\mathcal{R}_\omega$ under the \(\exp(-\jj\omega t)\) convention of
Chapter~\ref{ch:01}.

Let \(\Omega\subset\mathbb{R}^{3}\) be a piecewise smooth domain with boundary
\(\Gamma=\partial\Omega\). Phasor fields \(\E(\mathbf{r})\) and \(\H(\mathbf{r})\)
are suppressed in time unless causality is restored below. Impressed electric
current \(\Jimp\), charge density \(\rho\), and the data tuple \(\mathfrak{D}\) of
Eq.~\eqref{eq:ch1.data-specification-tuple} are not redefined; they are the
inputs on which the maps act. Finite-source admissibility is assumed in the sense of the
\ref{ch:01} $\mathcal{D}_J=1$ conditions
(Eq.~\eqref{eq:ch1.continuity-admissibility}): only configurations with
\(\mathcal{D}_{J}=1\) enter the domain of the radiation chain developed here
\cite{stratton1941}.

\begin{figure}[t]
\centering
\VolIPhysFig{\VolIPhysSourceRadiationSequence}
\caption{Excitation-to-observable chain: \(\mathcal{S}_\omega\) enforces Maxwell
closure with inherited boundary data; \(\Pi_{\mathrm{rad}}\) isolates the
radiative component on which far-zone flux integrals are meaningful.}
\label{fig:ch2.source-radiation-sequence}
\end{figure}

Function spaces are declared before the maps. Let $\mathcal{F}_\Omega$ be the
space of phasor pairs \((\E,\H)\) on \(\Omega\) whose components have
\(L^{2}\)-integrable curls and divergences in the sense required by
Eq.~\eqref{eq:ch1.maxwell-operator-block}, and whose boundary traces lie in the
domain of \(\mathcal{B}_\Gamma\). Let \(\mathcal{J}_\Omega\subset L^{2}(\Omega)^{3}\)
be the impressed-current space compatible with finite-energy class
(Section~\ref{sec:ch01.0.3}), charge continuity
(Eq.~\eqref{eq:ch1.charge-continuity-harmonic}), and the same boundary class.
Refined Sobolev and distributional extensions of these spaces are supplied in
Section~\ref{sec:ch211}; here the declaration is minimal but precise enough to
state the maps \cite{stratton1941}. The harmonic Maxwell \emph{solution map} is the
first quantitative object in the radiation chain: it records, for each admissible impressed
current \(\Jimp\in\mathcal{J}_\Omega\), the Maxwell field pair that the governing
equations and boundary data jointly permit. Physically, this map is what a full-wave solver
computes when it inverts Maxwell's equations on a declared domain; the output is a \emph{complete} field state, not yet a radiative report
\cite{stratton1941}. Experimentally one would observe, on a closed measurement surface
enclosing the source, both reactive circulation and outward-propagating flux in the
solution unless outgoing selection is applied afterward \cite{harrington2001}. The solution
map sends impressed data to a field state:
\begin{equation}
\label{eq:ch2.maxwell-solution-map}
\mathcal{S}_\omega:\mathcal{J}_\Omega\longrightarrow\mathcal{F}_\Omega,
\qquad
\Jimp\longmapsto(\E,\H),
\end{equation}
defined by \(\mathcal{S}_\omega[\Jimp]=(\E,\H)\) satisfying the Maxwell operator
residuals of Eq.~\eqref{eq:ch1.maxwell-operator-block} with sources \(\mathfrak{D}\)
and boundary data \(\mathbf{g}_\Gamma\) \cite{stratton1941}.
Equation~\eqref{eq:ch2.maxwell-solution-map} is the Chapter~\ref{ch:02} statement
of the initial-boundary value problem proven in Section~\ref{sec:ch01.1.1}; it
does not alter the law, only the input--output viewpoint.
Injectivity fails in general: distinct \(\Jimp\) can produce identical \((\E,\H)\)
on observational subdomains, and equivalent source distributions related by
divergence-free reparametrizations inside \(\Omega\) can leave exterior patterns
unchanged when \(\mathcal{B}_\Gamma\) permits reactive storage
(Section~\ref{sec:ch01.2.4}). Those equivalences are not nuisances to be removed
by definition; they foreshadow null-space and realizability analyses in
Sections~\ref{sec:ch27} and \ref{ch:04}. Radiation accounting therefore
never begins from a guessed modal field; it begins from \(\mathcal{S}_\omega[\Jimp]\)
with explicit completion.
The full state typically contains both stored and exported electromagnetic energy.
Reactive fields exchange power with sources and terminations at distances small
compared with \(k_0^{-1}\); radiative fields carry net outward flux at infinity
in open configurations. Let \(\mathcal{F}_{\mathrm{rad}}\subset\mathcal{F}_\Omega\)
denote pairs whose far-zone asymptotics satisfy outgoing selection (made analytic in
Section~\ref{sec:ch24}). Radiation accounting therefore requires a second map that
discards non-outgoing storage before pattern or transport claims are made. The outgoing
selector \(\Pi_{\mathrm{rad}}:\mathcal{F}_\Omega\to\mathcal{F}_{\mathrm{rad}}\) performs
that discard: it accepts a completed Maxwell pair and returns only the component whose
leading far-zone asymptotics satisfy Sommerfeld-outgoing qualification
\cite{stratton1941}. From a measurement standpoint, \(\Pi_{\mathrm{rad}}\) is the
mathematical counterpart of moving a probe beyond the reactive shell and reporting only
the transverse outgoing field that persists with \(1/r\) decay \cite{harrington2001}.
Composing source excitation with that selector defines the \emph{radiation operator}
\begin{equation}
\label{eq:ch2.radiation-operator-def}
\mathcal{R}_\omega:\mathcal{J}_\Omega\longrightarrow\mathcal{F}_{\mathrm{rad}},
\qquad
\mathcal{R}_\omega[\Jimp]=\Pi_{\mathrm{rad}}\,\mathcal{S}_\omega[\Jimp].
\end{equation}
Equation~\eqref{eq:ch2.radiation-operator-def} composes source excitation with
outgoing selection: \(\mathcal{R}_\omega\) accepts the same impressed currents as
\(\mathcal{S}_\omega\) but returns only the radiative image
\(\Pi_{\mathrm{rad}}(\E,\H)\in\mathcal{F}_{\mathrm{rad}}\), where
\(\mathcal{F}_{\mathrm{rad}}\) is the subspace of field pairs whose leading
far-zone behavior is outgoing and transverse.
The composition law records that radiation is never computed independently of
Maxwell completion:
\begin{equation}
\label{eq:ch2.operator-composition}
\mathcal{R}_\omega
=
\Pi_{\mathrm{rad}}\circ\mathcal{S}_\omega,
\qquad
\mathcal{D}_{J}=1
\Longrightarrow
\mathcal{S}_\omega:\mathcal{J}_\Omega\to\mathcal{F}_\Omega,
\quad
\mathcal{R}_\omega:\mathcal{J}_\Omega\to\mathcal{F}_{\mathrm{rad}}.
\end{equation}
where $\mathcal{S}_\omega$ is the harmonic Maxwell solution map, $\Pi_{\mathrm{rad}}$ the
outgoing selector, $\mathcal{R}_\omega$ the radiation map, $\mathcal{J}_\Omega$ the impressed-current
space, $\mathcal{F}_\Omega$ the admissible field space, $\mathcal{F}_{\mathrm{rad}}$ the radiative
subspace, and $\mathcal{D}_J=1$ the transport qualification inherited from
\ref{ch:01} \cite{stratton1941}. Equation~\eqref{eq:ch2.operator-composition} composes $\mathcal{S}_\omega$ with outgoing
selection $\Pi_{\mathrm{rad}}$ on $\mathcal{J}_\Omega$ \cite{stratton1941}. On a ball
$B_a$ with $k_0 a=0.5$, reactive-to-radiative field-energy ratios on $S_R$ with $R=2a$
typically exceed unity by an order of magnitude; applying pattern or transport metrics to
$\mathcal{S}_\omega[\Jimp]$ without $\Pi_{\mathrm{rad}}$ therefore overstates export by
roughly $10\times$ in magnitude bookkeeping before phase errors accumulate
\cite{jackson1999}. Superposition of radiative images is legitimate only after each
constituent satisfies Eq.~\eqref{eq:ch2.source-admissibility} and outgoing radiation
condition \cite{harrington2001}. Until Subsection~\ref{sec:ch224} fixes
$\Pi_{\mathrm{rad}}=\mathcal{O}_{\mathrm{out}}\circ\Pi_{\mathrm{prop}}$, treat
$\Pi_{\mathrm{rad}}$ as an abstract outgoing selector \cite{stratton1941}. Angular-flux
transport from Sections~\ref{sec:ch01.2.3}--\ref{sec:ch01.2.4} and structured-beam
interpretations in Section~\ref{sec:ch01.6.2} attach to $\mathcal{R}_\omega[\Jimp]$;
coefficients read from reactive components of $\mathcal{S}_\omega[\Jimp]$ describe
near-zone storage, not exported flux \cite{allen1992}.
When Green operators are available (Section~\ref{sec:ch23}), \(\mathcal{R}_\omega\)
admits a retarded integral representation that makes the propagation step explicit.
The integral representation below is the frequency-domain statement of what a
method-of-moments code evaluates when it integrates the dyadic Green function against
impressed currents: each source element at \(\mathbf{r}'\) contributes to the observation
point \(\mathbf{r}\) with phase and amplitude fixed by outgoing propagation
\cite{stratton1941}. Physically, the integral superposes the radiative
images of all current elements; only the portion of \(\mathcal{S}_\omega[\Jimp]\) that
\(\Pi_{\mathrm{rad}}\) retains enters the integrand in the completed theory developed in
Section~\ref{sec:ch23}. For \(\mathbf{r}\notin\Omega_s\) with
\(\operatorname{supp}\Jimp=\Omega_s\Subset\Omega\),
\begin{equation}
\label{eq:ch2.radiation-integral}
\E_{\mathrm{rad}}(\mathbf{r})
=
-\jj\omega\mu_0
\int_{\Omega_s}
\boldsymbol{\Gamma}(\mathbf{r},\mathbf{r}')\cdot\Jimp(\mathbf{r}')\,d^{3}r',
\end{equation}
In Eq.~\eqref{eq:ch2.radiation-integral},
\(\E_{\mathrm{rad}}(\mathbf{r})\) is the radiative electric phasor at
\(\mathbf{r}\) (V\,m\(^{-1}\));
\(\boldsymbol{\Gamma}(\mathbf{r},\mathbf{r}')\) is the outgoing dyadic Green tensor
(m\(^{-1}\)) that propagates unit current dipoles from \(\mathbf{r}'\) to
\(\mathbf{r}\);
\(\mu_0=4\pi\times 10^{-7}\,\mathrm{H\,m^{-1}}\) is the vacuum permeability;
\(\mathbf{r}'\) ranges over the source volume \(\Omega_s\); and the volume element
\(d^{3}r'\) carries SI units m\(^{3}\).
The factor \(-\jj\omega\mu_0\) arises from the harmonic Amp\`ere law under the
locked phasor convention.
Here \(\boldsymbol{\Gamma}\) is the outgoing dyadic Green function
\cite{stratton1941}. Equation~\eqref{eq:ch2.radiation-integral} is a
frequency-domain representation of the radiative field generated by a compact source
distribution through the outgoing dyadic Green function \cite{stratton1941}.
Every application of this representation assumes the $\mathcal{D}_J=1$ and boundary conditions
stated above. Unless noted otherwise, \(\omega>0\) is fixed,
materials are linear isotropic passive regions, and \(\Omega\) is either bounded
with declared \(\mathcal{B}_\Gamma\) or an exterior domain with compact
\(\operatorname{supp}\Jimp\) and outgoing radiation condition at infinity. Maps are defined
only on \(\mathcal{J}_\Omega\) satisfying Eq.~\eqref{eq:ch1.charge-continuity-harmonic};
generalized sources appear in Section~\ref{sec:ch232}. Maxwell linearity on \(\Omega\) induces linearity of both maps on fixed boundary
data: for \(\Jimp^{(1)},\Jimp^{(2)}\in\mathcal{J}_\Omega\) and scalars \(a_{1},a_{2}\),
\begin{equation}
\label{eq:ch2.radiation-linearity}
\mathcal{R}_\omega\!\left[a_{1}\Jimp^{(1)}+a_{2}\Jimp^{(2)}\right]
=
a_{1}\mathcal{R}_\omega[\Jimp^{(1)}]
+a_{2}\mathcal{R}_\omega[\Jimp^{(2)}],
\end{equation}
In Eq.~\eqref{eq:ch2.radiation-linearity},
\(a_1,a_2\in\mathbb{C}\) are complex amplitudes and
\(\Jimp^{(1)},\Jimp^{(2)}\in\mathcal{J}_\Omega\) are admissible impressed currents;
\(\mathcal{R}_\omega\) is the radiation map; and linearity holds when boundary data
\(\mathbf{g}_\Gamma\) and the material layout are held fixed \cite{maxwell1873}.
Composite-source analysis decomposes \(\Jimp\) into constituents,
maps each through \(\mathcal{R}_\omega\), and superposes radiative outputs---a
discipline that prevents reactive interference from masquerading as transport
(Section~\ref{sec:ch01.6.4}).
Write \(\operatorname{ran}\mathcal{S}_\omega\subset\mathcal{F}_\Omega\) for all
field pairs reachable from finite sources and
\(\operatorname{ran}\mathcal{R}_\omega\subset\mathcal{F}_{\mathrm{rad}}\) for
far-zone radiative outputs. Strict inclusion
\(\operatorname{ran}\mathcal{R}_\omega\subsetneq\operatorname{ran}\mathcal{S}_\omega\)
occurs whenever reactive storage is nontrivial: distinct near-zone states can share
identical radiative images (Section~\ref{sec:ch01.2.4}). Pattern measurements
probe \(\operatorname{ran}\mathcal{R}_\omega\) and cannot invert
\(\mathcal{S}_\omega\) without additional near-zone data---a limitation revisited in
Section~\ref{sec:ch244}.
Radiative power is metered with inherited Poynting machinery.
The practical question is how much time-averaged power crosses a closed surface that
encloses the source and lies beyond the reactive shell; that flux is the scalar
quantity laboratories report as radiated power.
On a sphere
\(S_R=\{|\mathbf{r}|=R\}\) enclosing \(\operatorname{supp}\Jimp\), with outward
unit vector \(\hat{\mathbf{r}}\) and complex Poynting vector
\(\mathbf{S}=\tfrac{1}{2}\E\times\H^*\) from Section~\ref{sec:ch01.1.4},
\begin{equation}
\label{eq:ch2.radiated-power-screen}
P_{\mathrm{rad}}(R)
=
\Re\oint_{S_R}\mathbf{S}\cdot\hat{\mathbf{r}}\,dS.
\end{equation}
In Eq.~\eqref{eq:ch2.radiated-power-screen},
\(P_{\mathrm{rad}}(R)\) is the time-averaged outward power (W) through the sphere
of radius \(R\) (m);
\(\mathbf{S}\) is the complex Poynting vector (W\,m\(^{-2}\)) built from the
radiative field pair;
\(\hat{\mathbf{r}}=\mathbf{r}/|\mathbf{r}|\) is the outward unit normal on
\(S_R\); and \(dS\) is the area element (m\(^{2}\)).
The real part \(\Re\{\cdot\}\) extracts the time-averaged power flux carried by the
phasor fields under the \(\exp(-\jj\omega t)\) convention.
Equation~\eqref{eq:ch2.radiated-power-screen} inherits the Poynting law of
Section~\ref{sec:ch01.1.4}. For \(\mathcal{R}_\omega[\Jimp]\), \(P_{\mathrm{rad}}(R)\) becomes
$R$-independent beyond the reactive shell \cite{poynting1884}. Angular-momentum budgets from
Section~\ref{sec:ch01.1.5} attach to the same radiative component.
\begin{figure}[t]
\centering
\ChOneFigBlock{%
\begin{tikzpicture}[ch1fig, node distance=5mm and 7mm]
 \node[ch1box, minimum width=22mm] (core) at (0,0) {Source\\$\Omega_s$};
 \node[ch1box, minimum width=28mm] (react) at (0,-12mm) {Reactive shell\\$1/r^2$, $1/r^3$};
 \node[ch1box, minimum width=32mm] (sel) at (0,-24mm) {Selector\\$\Pi_{\mathrm{rad}}$};
 \node[ch1box, minimum width=28mm] (far) at (0,-36mm) {Far zone\\$1/r$ outgoing};
 \draw[ch1flow] (core) -- (react) -- (sel) -- (far);
\end{tikzpicture}%
}
\caption{Spatial nesting seen by the maps: \(\mathcal{S}_\omega\) populates all
regions; \(\Pi_{\mathrm{rad}}\) passes only the outgoing transverse content
compatible with Section~\ref{sec:ch01.2.4}.}
\label{fig:ch2.reactive-radiative-split}
\end{figure}
The solution map energizes reactive and radiative regions alike; the
radiation map discards non-outgoing storage before pattern or transport claims are
made. Misidentifying reactive circulation as exported angular structure is a
recurring error in high-gradient near-zone beams \cite{allen1992}. Finite measurement channels (Section~\ref{sec:ch01.5}) implement linear functionals
on field space. Composing a declared far-zone observable with \(\mathcal{R}_\omega\)
yields a linear functional of \(\Jimp\); composing the same observable with
\(\mathcal{S}_\omega\) without selection generally includes reactive contamination.
No new measurement axioms are introduced---only the correct map for the declared
observable class \cite{stratton1941}. Non-radiating currents illustrate kernel visibility. If
\(\Jimp^{\mathrm{(NR)}}\in\mathcal{J}_\Omega\) has vanishing far-zone transverse
fields while maintaining near-zone storage (Section~\ref{sec:ch01.2.4}), then
\begin{equation}
\label{eq:ch2.null-radiation}
\mathcal{R}_\omega[\Jimp^{\mathrm{(NR)}}]=\mathbf{0}
\quad\text{in}\ \mathcal{F}_{\mathrm{rad}},
\end{equation}
In Eq.~\eqref{eq:ch2.null-radiation},
\(\Jimp^{\mathrm{(NR)}}\) is an admissible impressed current (A\,m\(^{-2}\)) that satisfies
\(\mathcal{D}_{J}=1\) yet carries no outgoing transverse export;
\(\mathcal{F}_{\mathrm{rad}}\) is the far-zone radiative field space selected by
\(\Pi_{\mathrm{rad}}\); and \(\ker\mathcal{R}_\omega\) denotes the set of currents for
which the radiation map vanishes on that space,
while \(\mathcal{S}_\omega[\Jimp^{\mathrm{(NR)}}]\neq\mathbf{0}\) in general.
Equation~\eqref{eq:ch2.null-radiation} targets the null space analyzed in
Section~\ref{sec:ch272}; ignoring \(\ker\mathcal{R}_\omega\) invites spurious modal
assignments to invisible sources.
Membership in \(\mathcal{J}_\Omega\) requires more than square integrability:
the Maxwell-closed screen below states, in compact form, which impressed currents
may enter the radiation chain without contradicting charge continuity or boundary
compatibility. \(\mathcal{D}_{J}\) is not a mathematical
regularity condition alone: it encodes the laboratory requirement that impressed currents
be charge-consistent and boundary-compatible before Maxwell field specification is attempted
\cite{maxwell1873}. Violating the screen produces formal solutions that no
passive feed can realize \cite{harrington2001}:
\begin{equation}
\label{eq:ch2.source-admissibility}
\mathcal{D}_{J}[\Jimp]=1
\iff
\left(
\Jimp\in L^{2}(\Omega)^{3},\
\divg\Jimp-\jj\omega\rho=0,\
\Jimp\ \text{satisfies BC compatibility on }\Gamma
\right),
\end{equation}
In Eq.~\eqref{eq:ch2.source-admissibility},
\(\mathcal{D}_{J}\) is the source Maxwell-closed screen (dimensionless) inherited from
Eq.~\eqref{eq:ch1.charge-continuity-harmonic};
\(\Jimp\) is the impressed electric current density (A\,m\(^{-2}\)) on compact
\(\Omega\subset\mathbb{R}^{3}\);
\(\rho\) is the impressed charge density (C\,m\(^{-3}\)) linked by
\(\divg\Jimp-\jj\omega\rho=0\); and \(\Gamma=\partial\Omega\) is the boundary on which
tangential traces must match the declared boundary class.
with \(\mathcal{D}_{J}\) as in Eq.~\eqref{eq:ch1.continuity-admissibility}.
Violating \(\divg\Jimp-\jj\omega\rho=0\) contradicts
Eq.~\eqref{eq:ch1.charge-continuity-harmonic} before Maxwell field specification is attempted.
Boundary compatibility ensures induced surface equivalents do not contradict
Eq.~\eqref{eq:ch1.boundary-data-operator} \cite{stratton1941}.
Causality constrains harmonic representations.
A frequency-domain phasor is physically meaningful only when it can be traced to a
retarded time-domain evolution that respects finite propagation speed.
With time-domain
current \(j(\mathbf{r},t)\) and retarded dyadic impulse
\(\mathcal{G}(\mathbf{r},\mathbf{r}',t-t')\),
\begin{equation}
\label{eq:ch2.causality-retarded}
\E(\mathbf{r},t)
=
-\mu_0\,\frac{\partial}{\partial t}
\int
\mathcal{G}(\mathbf{r},\mathbf{r}',t-t')\cdot j(\mathbf{r}',t')\,d^{3}r'\,dt',
\end{equation}
integrated only over \(t'\le t\) \cite{stratton1941}.
In Eq.~\eqref{eq:ch2.causality-retarded},
\(\E(\mathbf{r},t)\) is the real time-domain electric field (V\,m\(^{-1}\));
\(j(\mathbf{r}',t')\) is the impressed current density (A\,m\(^{-2}\)) at source
point \(\mathbf{r}'\) and retarded time \(t'\);
\(\mathcal{G}(\mathbf{r},\mathbf{r}',t-t')\) is the causal dyadic Green impulse
response (m\(^{-1}\)\,s\(^{-1}\)) that vanishes for \(t-t'<0\); and the temporal
integral enforces that the field at \((\mathbf{r},t)\) depends only on sources in
the past light cone.
The steady-state operator
$\mathcal{R}_\omega$ may therefore be viewed as the frequency-domain specialization of
the retarded causal evolution \cite{stratton1941}. The equivalence between
the time-domain and harmonic descriptions follows under the standard assumptions of
linearity, stationarity, and finite-energy excitation \cite{maxwell1873}.
Vacuum scales from Eq.~\eqref{eq:ch1.vacuum-impedance-wavenumber} fix the carrier
on which selection acts: \(k_0=\omega\sqrt{\mu_0\varepsilon_0}\),
\(\eta_0=\sqrt{\mu_0/\varepsilon_0}\). Outgoing components retained by
\(\Pi_{\mathrm{rad}}\) carry \(\exp(-\jj k_0 r)\) phase fronts at leading order;
components whose phase does not lock to outgoing \(k_0 r\) remain in
\(\mathcal{S}_\omega[\Jimp]\) but outside \(\mathcal{F}_{\mathrm{rad}}\). Power
metered by Eq.~\eqref{eq:ch2.radiated-power-screen} stabilizes only on the latter
\cite{stratton1941}. Distance labels become meaningful through operator action. For compact
\(\operatorname{supp}\Jimp\), \(\mathcal{S}_\omega[\Jimp]\) contains
\(1/r^{2}\) and \(1/r^{3}\) reactive terms; \(\mathcal{R}_\omega[\Jimp]\) retains
leading \(1/r\) transverse fields with \(kr\) phase in the far zone
(Section~\ref{sec:ch01.2.4}). Asymptotic partitions are sharpened in
Section~\ref{sec:ch252}.
Observability enters as functional composition rather than as a new postulate.
Let \(\mathcal{O}_{\mathrm{far}}\) denote a linear functional that extracts a
declared far-zone scalar or vector observable---pattern voltage, axial ratio
proxy, or a single spherical-harmonic coefficient---from a field pair in
\(\mathcal{F}_{\mathrm{rad}}\). The measured quantity associated with impressed current \(\Jimp\) is the scalar
\(\mathcal{O}_{\mathrm{far}}\big[\mathcal{R}_\omega[\Jimp]\big]\), not
\(\mathcal{O}_{\mathrm{far}}\big[\mathcal{S}_\omega[\Jimp]\big]\) unless the
observable has been shown to vanish on reactive content. Because both
\(\mathcal{R}_\omega\) and \(\mathcal{O}_{\mathrm{far}}\) are linear on their
domains, the laboratory reading is a linear functional
\(\Lambda_{\mathrm{rad}}=\mathcal{O}_{\mathrm{far}}\circ\mathcal{R}_\omega\)
on \(\mathcal{J}_\Omega\). That composition expresses in map language
the measurement limitations in Section~\ref{sec:ch01.5}: finite apertures and
finite modal sets appear as rank-deficient \(\mathcal{O}_{\mathrm{far}}\), not
as modifications of Maxwell law \cite{stratton1941}. Angular-transport claims
require further that \(\Lambda_{\mathrm{rad}}\) be compatible with the flux
budgets of Section~\ref{sec:ch01.1.5} when the observable is interpreted as
momentum or power export.
The outgoing selector \(\Pi_{\mathrm{rad}}\) is idempotent on
\(\mathcal{F}_{\mathrm{rad}}\): applying selection twice leaves a radiative state
unchanged. On the larger space \(\mathcal{F}_\Omega\), \(\Pi_{\mathrm{rad}}\) is
a projection only after the Sommerfeld-class uniqueness of Section~\ref{sec:ch243}
is imposed; without that completion, multiple outgoing decompositions can coexist
for the same truncated near-zone data. The idempotence condition is therefore a
\emph{joint} property of field space and boundary class, not a algebraic trick
independent of physics \cite{harrington2001}. Numerical methods that
window fields in space and call the result ``radiative'' without outgoing radiation condition are approximating \(\Pi_{\mathrm{rad}}\) by a non-idempotent mask; such outputs
must not be fed into transport criteria without an explicit error budget.
The composition law used downstream is
\begin{equation}
\label{eq:ch2.measurement chain-composition}
\mathcal{R}_\omega
=
\Pi_{\mathrm{rad}}\circ\mathcal{S}_\omega,
\qquad
\mathcal{D}_{J}=1
\Longrightarrow
\mathcal{S}_\omega:\mathcal{J}_\Omega\to\mathcal{F}_\Omega,
\quad
\mathcal{R}_\omega:\mathcal{J}_\Omega\to\mathcal{F}_{\mathrm{rad}}.
\end{equation}
In Eq.~\eqref{eq:ch2.measurement chain-composition},
\(\mathcal{S}_\omega\) is the Maxwell solution map from admissible sources to field pairs,
\(\Pi_{\mathrm{rad}}\) is the outgoing selector that retains far-zone transverse export,
\(\mathcal{R}_\omega=\Pi_{\mathrm{rad}}\circ\mathcal{S}_\omega\) is the radiation map,
\(\mathcal{D}_{J}=1\) denotes source Maxwell-closed from
Eq.~\eqref{eq:ch2.source-admissibility}, and \(\mathcal{J}_\Omega\),
\(\mathcal{F}_\Omega\), \(\mathcal{F}_{\mathrm{rad}}\) are the source and field spaces
declared in Subsection~\ref{sec:ch201}.
Green kernels realize \(\mathcal{S}_\omega\); Sommerfeld conditions define
\(\Pi_{\mathrm{rad}}\); spectral theory in Section~\ref{sec:ch26} studies
eigenstructure on spaces refined in Section~\ref{sec:ch21}
\cite{harrington2001}. Nothing in
Eqs.~\eqref{eq:ch2.maxwell-solution-map}--\eqref{eq:ch2.measurement chain-composition}
re-proves results listed in Section~\ref{sec:ch01.7.1}; continuity
Eq.~\eqref{eq:ch1.continuity-map} is honored by citation-only continuity
of curl--curl algebra, stress flux theorems, and gauge lemmas.
Subsection~\ref{sec:ch201} has named \(\mathcal{S}_\omega\), \(\Pi_{\mathrm{rad}}\), and
\(\mathcal{R}_\omega\), introduced their Green and Poynting realizations, and connected
finite measurement rank to composition with far-zone observables. Subsection~\ref{sec:ch202}
identifies the Maxwell operator whose inversion, under declared boundary completion,
produces \(\mathcal{S}_\omega\).
